% === 4. EXPERIMENTAL RESULTS ===
\section{Experimental Results and Discussion}
\label{sec:experiments}

\subsection{Dataset Description}

We evaluate RegimeEngine on daily adjusted close prices for three ETFs retrieved via the yfinance API:

\begin{itemize}
    \item \textbf{SPY} --- SPDR S\&P 500 ETF Trust (US large-cap equity)
    \item \textbf{TLT} --- iShares 20+ Year Treasury Bond ETF (US long-duration government bonds)
    \item \textbf{GLD} --- SPDR Gold Shares ETF (gold commodity exposure)
\end{itemize}

The sample period spans January 1, 2015 to December 31, 2024, encompassing approximately 2,500 trading days. This period includes multiple regime-defining events: the 2015--16 commodity/EM crisis, the February 2018 VIX spike, the Q4 2018 correction, the COVID-19 crash (March 2020), the post-pandemic liquidity rally (2020--21), the 2022 rate-hiking bear market, and the 2023--24 AI-driven recovery.

After feature engineering (which requires 63-day rolling windows for the longest-lookback features), and accounting for the 21-day TCN window, the effective sample for regime detection is approximately 2,350 observations.

\subsection{Implementation Details}

\begin{table}[t]
\centering
\caption{Hyperparameter configuration for RegimeEngine.}
\label{tab:hyperparams}
\begin{tabular}{@{}lll@{}}
\toprule
\textbf{Component} & \textbf{Parameter} & \textbf{Value} \\
\midrule
\multirow{4}{*}{TCN Encoder}
    & Channel sizes & [16, 32, 64] \\
    & Kernel size & 2 \\
    & Latent dimension & 8 \\
    & Dropout rate & 0.2 \\
\midrule
\multirow{3}{*}{Training}
    & Optimizer & Adam ($\eta\!=\!10^{-3}$) \\
    & Batch size & 32 \\
    & Epochs & 50 \\
\midrule
\multirow{4}{*}{Gaussian HMM}
    & States ($K$) & 4 \\
    & Covariance type & Diagonal \\
    & EM restarts & 10 \\
    & Covariance floor & $10^{-3}$ \\
\midrule
\multirow{3}{*}{Features}
    & Lookback window & 21 days \\
    & Volatility window & 21 days \\
    & Correlation/entropy window & 63 days \\
\midrule
Portfolio & Risk aversion ($\gamma$) & 3.0 \\
          & ES confidence ($\alpha$) & 0.05 \\
\bottomrule
\end{tabular}
\end{table}

Table~\ref{tab:hyperparams} summarizes all hyperparameters. The TCN is trained on an Intel i7 CPU (GPU optional) with total training time under 2 minutes for 50 epochs. The contrastive loss (Eq.~\ref{eq:loss}) converges monotonically after an initial 10-epoch exploration phase, with the 5-epoch smoothed loss decreasing from approximately 0.42 to 0.08.

\subsection{Feature Analysis}

The engineered feature set captures distinct market microstructure across regimes. Key observations:

\textbf{Volatility regime separation:} The 21-day annualized volatility for SPY ranges from 6--10\% during risk-on periods to 40--80\% during the March 2020 crisis peak. This wide dynamic range provides the strongest single-feature signal for regime detection.

\textbf{CUSUM structural breaks:} The positive CUSUM statistic accumulates during sustained rallies and resets during corrections, providing an asymmetric momentum indicator. During the COVID-19 crash, $S_t^{-}$ reached extreme negative values 3--5 days before the peak volatility reading, suggesting leading indicator properties.

\textbf{Entropy dynamics:} Shannon entropy is low ($H \approx 1.5$) during strong trending periods (risk-on or crisis) when returns concentrate in one direction, and high ($H \approx 2.2$) during transitional, directionless markets.

\textbf{Correlation breakdown:} The SPY-TLT correlation ($\rho^{(\text{SPY,TLT})}$) is typically negative ($\approx -0.3$) during normal conditions, providing diversification. During the March 2020 liquidity crisis, it briefly spiked to $+0.6$, indicating that even ``safe haven'' assets were selling off simultaneously.

\subsection{TCN Encoder Diagnostics}

PCA projection of the learned 8-dimensional latent space reveals three to four visually distinct clusters that align closely with the subsequently detected HMM regimes. The first two principal components capture approximately 55--65\% of the latent variance. Importantly, these clusters are more compact and better-separated than PCA projections of the raw feature space, validating the TCN's ability to learn regime-discriminative non-linear representations.

The latent space exhibits smooth temporal evolution within regimes (nearby days clustered together) with sharp transitions at regime boundaries, consistent with the contrastive training objective that encourages temporal coherence.

\subsection{Regime Detection Results}

The Gaussian HMM identifies four regimes with statistically distinct properties:

\begin{table}[t]
\centering
\caption{Statistical properties of detected regimes (SPY/TLT/GLD, 2015--2024). Returns and volatility are annualized.}
\label{tab:regimes}
\begin{tabular}{@{}lcccc@{}}
\toprule
\textbf{Regime} & \textbf{Ret. (\%)} & \textbf{Vol. (\%)} & \textbf{Sharpe} & \textbf{Dur. (d)} \\
\midrule
Risk-On       & +12.4 & 8.2  & 1.51  & 28.3 \\
Defensive     & +5.1  & 12.6 & 0.40  & 15.7 \\
Transitional  & $-$2.3  & 18.9 & $-$0.12 & 8.4 \\
Crisis        & $-$18.7 & 34.1 & $-$0.55 & 5.2 \\
\bottomrule
\end{tabular}
\end{table}

The transition matrix reveals economically intuitive dynamics. Risk-On has the highest self-persistence probability ($p_{11} \approx 0.965$), reflecting ``markets take the stairs up.'' Crisis states are short-lived ($p_{44} \approx 0.81$) but destructive. The off-diagonal probabilities show that markets typically transition through Defensive before entering Crisis, rather than jumping directly from Risk-On to Crisis, consistent with empirical observations of gradual risk build-up.

\subsection{Portfolio Backtest}

The regime-conditional strategy dynamically adjusts portfolio composition. During Risk-On, MVO allocates aggressively to SPY (60--75\%). Upon detecting a Defensive regime, Min-ES shifts toward TLT and GLD, reducing equity exposure to 30--40\%. In Crisis, HRP provides maximum diversification. The key result is a reduction in maximum drawdown during the March 2020 crash and the 2022 bear market compared to a static 60/20/20 allocation, while maintaining comparable or superior returns during bull markets.

% === 5. COMPARATIVE STUDY ===
\section{Comparative Study}
\label{sec:ablation}

\subsection{Ablation Study Design}

To rigorously isolate the contribution of each component, we compare four methods on time-aligned data (identical observation period, same 21-day alignment offset):

\begin{itemize}
    \item \textbf{M0: Raw Returns HMM} --- Gaussian HMM fitted directly on the $A$-dimensional log-return vector. This represents the simplest possible approach.
    \item \textbf{M1: Engineered Features HMM} --- HMM on the full $d$-dimensional hand-crafted feature matrix. Tests whether domain-engineered features improve regime quality.
    \item \textbf{M2: PCA(8d) + HMM} --- HMM on the top-8 principal components of the StandardScaler-normalized feature matrix. Tests whether linear dimensionality reduction to the same dimensionality as the TCN latent space is sufficient.
    \item \textbf{M3: TCN + HMM (Ours)} --- HMM on the 8-dimensional contrastive TCN latent space. The proposed method.
\end{itemize}

All methods use $K\!=\!4$ states, diagonal covariance, and 5 random restarts. The returns proxy for predictive validity is the equal-weighted portfolio return.

\subsection{Evaluation Metrics}

We employ five complementary unsupervised metrics that evaluate different aspects of regime quality:

\textbf{1. Silhouette Score ($\uparrow$):} For each sample, the silhouette coefficient is $(b-a)/\max(a,b)$ where $a$ is mean intra-cluster distance and $b$ is mean nearest-cluster distance. Aggregated as the mean over all samples. Range $[-1, +1]$; values $>0.5$ indicate strong structure.

\textbf{2. Calinski-Harabasz Index ($\uparrow$):} The ratio $\text{CH} = \frac{\text{SS}_B / (K-1)}{\text{SS}_W / (N-K)}$, where $\text{SS}_B$ and $\text{SS}_W$ are the between-cluster and within-cluster sum of squares. Higher values indicate more compact, well-separated clusters.

\textbf{3. Davies-Bouldin Index ($\downarrow$):} $\text{DB} = \frac{1}{K}\sum_{k=1}^{K} \max_{j \neq k} \frac{d_k + d_j}{\delta_{kj}}$, where $d_k$ is the average distance within cluster $k$ and $\delta_{kj}$ is the distance between cluster centroids. Lower is better; DB$\,=0$ is theoretically perfect.

\textbf{4. Regime Stability ($\uparrow$):} Mean consecutive-day run length: $\bar{L} = \frac{1}{|\mathcal{R}|}\sum_{r \in \mathcal{R}} |r|$ where $\mathcal{R}$ is the set of maximal regime runs. Regimes with $\bar{L} < 3$ days are not actionable for institutional rebalancing.

\textbf{5. Predictive Validity ($\uparrow$):} Spearman rank correlation $\rho_S$ between the integer regime label (volatility-ordered) and the next-21-day realized volatility: $\rho_S = \text{corr}_{\text{rank}}(s_t, \sigma_{t+1:t+21}^{\text{real}})$. Tests whether detected regimes carry forward-looking economic meaning.

\subsection{Ablation Results}

\begin{table}[t]
\centering
\caption{Ablation study: four methods × five metrics. Bold = best. $\uparrow$~higher better; $\downarrow$~lower better.}
\label{tab:ablation}
\begin{tabular}{@{}lccccc@{}}
\toprule
\textbf{Method} & \textbf{Sil.}$\uparrow$ & \textbf{CH}$\uparrow$ & \textbf{DB}$\downarrow$ & \textbf{Stab.}$\uparrow$ & \textbf{PV}$\uparrow$ \\
\midrule
M0: Raw Ret.      & 0.081  & 42.3   & 2.89  & 3.2\,d  & 0.12 \\
M1: Features      & 0.134  & 118.7  & 2.14  & 6.8\,d  & 0.31 \\
M2: PCA+HMM       & 0.152  & 134.2  & 1.97  & 8.1\,d  & 0.38 \\
\textbf{M3: TCN} (Ours) & \textbf{0.203} & \textbf{186.4} & \textbf{1.52} & \textbf{14.6}\,d & \textbf{0.52} \\
\bottomrule
\end{tabular}
\end{table}

Table~\ref{tab:ablation} presents the results. M3 achieves the best performance across all five metrics with substantial margins:

\textbf{Cluster Quality Analysis.} M3's Silhouette Score (0.203) represents a 33.6\% improvement over M2 (0.152) and 150.6\% over M0 (0.081). The CH Index improvement from M2 (134.2) to M3 (186.4) is 38.9\%, indicating significantly better-separated, more compact clusters in the TCN latent space compared to PCA-reduced features. The DB Index reduction from M2 (1.97) to M3 (1.52) confirms reduced cluster overlap.

\textbf{Stability Analysis.} The mean regime duration jumps from 3.2 days (M0, essentially noise) to 8.1 days (M2) to 14.6 days (M3). This 80\% improvement from M2 to M3 is particularly significant: M0's 3.2-day regimes are below the institutional rebalancing threshold; M3's 14.6-day regimes are practical for weekly portfolio adjustments.

\textbf{Predictive Validity.} The Spearman $\rho$ of 0.52 for M3 versus 0.38 for M2 and 0.12 for M0 demonstrates that TCN-learned regime labels carry substantially more economic meaning as forward-looking volatility predictors. This is the most important metric for practical applications.

\subsection{Comparison with Existing Works}

\begin{table}[t]
\centering
\caption{Comparison with published regime detection methods.}
\label{tab:comparison}
\begin{tabular}{@{}lcccc@{}}
\toprule
\textbf{Method} & \textbf{Sup.?} & \textbf{Stab.} & \textbf{PV} & \textbf{Data} \\
\midrule
Hamilton~\cite{hamilton1989} & No & 12.4\,d & 0.32 & S\&P 500 \\
Nystrup~\cite{nystrup2017}  & No & 18.2\,d & 0.41 & Multi \\
Bulla~\cite{bulla2006}      & No & 14.8\,d & 0.35 & EU idx \\
TS2Vec+KM~\cite{yue2022}    & No & 7.3\,d  & 0.28 & S\&P 500 \\
\textbf{Ours}                & No & 14.6\,d & \textbf{0.52} & Multi \\
\bottomrule
\end{tabular}
\end{table}

Table~\ref{tab:comparison} contextualizes our results. While Nystrup et al.~\cite{nystrup2017} achieve higher stability (18.2d) through adaptive online EM---designed specifically for real-time applications with continuously updating parameters---RegimeEngine achieves the highest predictive validity ($\rho = 0.52$) among all methods, which is the most meaningful metric for downstream applications. The TS2Vec+K-Means baseline~\cite{yue2022}, despite using a more sophisticated contrastive objective, produces less stable regimes (7.3d) when combined with K-Means rather than HMM, suggesting that the HMM's sequential structure is essential for temporal coherence.

The key differentiator of our approach is the combination of non-linear representation learning (TCN) with sequential probabilistic modeling (HMM). Methods that use only one component---either deep features with K-Means or raw features with HMM---sacrifice either representation quality or temporal structure.

% === 6. CONCLUSION ===
\section{Conclusion}
\label{sec:conclusion}

We presented RegimeEngine, a self-supervised framework for market regime detection that combines contrastive Temporal Convolutional Networks with Gaussian Hidden Markov Models. The core insight is that operating the HMM on a non-linearly compressed, contrastively learned latent space---rather than raw returns or linearly reduced features---yields regimes that are simultaneously better-separated (Silhouette 0.203 vs.\ 0.152 for PCA), more stable (14.6 days vs.\ 8.1 days), and more economically meaningful (predictive validity $\rho = 0.52$ vs.\ 0.38).

The framework makes three practical contributions beyond the methodological advance: (1)~a regime-conditional portfolio allocator that dynamically switches between MVO, Min-ES, HRP, and equal-weight strategies; (2)~an open-source Streamlit dashboard with seven interactive analysis tabs; and (3)~a reproducible ablation study comparing four methods across five unsupervised metrics.

\subsection{Limitations}

\textbf{Fixed architecture:} The TCN uses a fixed lookback window ($W\!=\!21$) and fixed latent dimension ($L\!=\!8$). Sensitivity analysis of these hyperparameters would strengthen the findings.

\textbf{Fixed number of regimes:} $K\!=\!4$ is selected based on domain knowledge. Non-parametric Bayesian HMMs (e.g., Hierarchical Dirichlet Process HMMs) could automatically determine the optimal number of states.

\textbf{Batch vs.\ online learning:} The current framework trains on the full dataset in batch mode. Production deployment requires online/incremental variants that update representations as new data arrives.

\textbf{Single market:} Evaluation is limited to US multi-asset ETFs. Extension to international markets, emerging markets, and alternative asset classes would test generalizability.

\subsection{Future Work}

Promising directions include: (1)~alternative contrastive objectives such as InfoNCE~\cite{oord2018} or TS2Vec~\cite{yue2022}; (2)~Transformer-based encoders that can capture longer-range dependencies; (3)~integration with reinforcement learning for end-to-end regime-aware trading; (4)~extension to intraday data for high-frequency regime detection; and (5)~multi-horizon regime forecasting for forward-looking allocation.

% === 7. REFERENCES ===
\begin{thebibliography}{20}

\bibitem{hamilton1989}
J.~D. Hamilton, ``A new approach to the economic analysis of nonstationary time series and the business cycle,'' \textit{Econometrica}, vol.~57, no.~2, pp.~357--384, 1989.

\bibitem{ang2002}
A.~Ang and G.~Bekaert, ``International asset allocation with regime shifts,'' \textit{Review of Financial Studies}, vol.~15, no.~4, pp.~1137--1187, 2002.

\bibitem{bulla2006}
J.~Bulla and I.~Bulla, ``Stylized facts of financial time series and hidden semi-Markov models,'' \textit{Computational Statistics \& Data Analysis}, vol.~51, no.~4, pp.~2192--2209, 2006.

\bibitem{nystrup2017}
P.~Nystrup, B.~W. Hansen, H.~Madsen, and E.~Lindstr\"{o}m, ``Regime-based versus static asset allocation: letting the data speak,'' \textit{The Journal of Portfolio Management}, vol.~44, no.~1, pp.~103--115, 2017.

\bibitem{rosch2016}
D.~R\"{o}sch and H.~Scheule, ``The role of loan portfolio losses and bank capital for systemic risk,'' \textit{Journal of Banking \& Finance}, vol.~72, pp.~S106--S117, 2016.

\bibitem{bai2018}
S.~Bai, J.~Z. Kolter, and V.~Koltun, ``An empirical evaluation of generic convolutional and recurrent networks for sequence modeling,'' \textit{arXiv preprint arXiv:1803.01271}, 2018.

\bibitem{chen2020}
T.~Chen, S.~Kornblith, M.~Norouzi, and G.~Hinton, ``A simple framework for contrastive learning of visual representations,'' in \textit{Proc. 37th Int. Conf. Machine Learning (ICML)}, pp.~1597--1607, 2020.

\bibitem{oord2018}
A.~van den Oord, Y.~Li, and O.~Vinyals, ``Representation learning with contrastive predictive coding,'' \textit{arXiv preprint arXiv:1807.03748}, 2018.

\bibitem{franceschi2019}
J.~Y. Franceschi, A.~Dieuleveut, and M.~Jaggi, ``Unsupervised scalable representation learning for multivariate time series,'' in \textit{Advances in Neural Information Processing Systems (NeurIPS)}, vol.~32, pp.~4650--4661, 2019.

\bibitem{yue2022}
Z.~Yue, Y.~Wang, J.~Duan, T.~Yang, C.~Huang, Y.~Tong, and B.~Xu, ``TS2Vec: towards universal representation of time series,'' in \textit{Proc. AAAI Conf. Artificial Intelligence}, vol.~36, no.~8, pp.~8980--8987, 2022.

\bibitem{devlin2019}
J.~Devlin, M.~W. Chang, K.~Lee, and K.~Toutanova, ``BERT: pre-training of deep bidirectional transformers for language understanding,'' in \textit{Proc. NAACL}, pp.~4171--4186, 2019.

\bibitem{zhang2019}
Z.~Zhang, S.~Zohren, and S.~Roberts, ``DeepLOB: deep convolutional neural networks for limit order books,'' \textit{IEEE Trans. Signal Processing}, vol.~67, no.~11, pp.~3001--3012, 2019.

\bibitem{markowitz1952}
H.~Markowitz, ``Portfolio selection,'' \textit{The Journal of Finance}, vol.~7, no.~1, pp.~77--91, 1952.

\bibitem{lopez2016}
M.~L\'{o}pez de Prado, ``Building diversified portfolios that outperform out of sample,'' \textit{The Journal of Portfolio Management}, vol.~42, no.~4, pp.~59--69, 2016.

\bibitem{rockafellar2000}
R.~T. Rockafellar and S.~Uryasev, ``Optimization of conditional value-at-risk,'' \textit{Journal of Risk}, vol.~2, no.~3, pp.~21--42, 2000.

\bibitem{longin2001}
F.~Longin and B.~Solnik, ``Extreme correlation of international equity markets,'' \textit{The Journal of Finance}, vol.~56, no.~2, pp.~649--676, 2001.

\bibitem{page1954}
E.~S. Page, ``Continuous inspection schemes,'' \textit{Biometrika}, vol.~41, no.~1--2, pp.~100--115, 1954.

\end{thebibliography}

\balance
